\subsection{WEYL GROUP}

Let G be a connected compact Lie group and T a maximal torus of G. Denote by \(N_{G}(T)\) the normalizer of T in G; by Prop. 4 (no. 4), the quotient group \(N_{G}(T)/T\) is finite. We denote it by \(W_{G}(T)\) , or by \(W(T)\) , and call it the Weyl group of the maximal torus T of G, or the Weyl group of G relative to T. Since T is commutative, the operation of \(N_{G}(T)\) on T by inner automorphisms of G induces by passage to the quotient an operation, called the canonical operation, of the group \(W_{G}(T)\) on the Lie group T. By Cor. 6 of Th. 2 of no. 2, this operation is faithful: the associated homomorphism \(W_{G}(T) \rightarrow Aut T\) is injective.

If \(T'\) is another maximal torus of G and if \(g \in G\) is such that Int g maps T to \(T'\) (no. 2, Th. 2 b)), then Int g induces an isomorphism \(a_g\) from \(\mathrm{W}_{\mathrm{G}}(\mathrm{T})\) to \(\mathrm{W}_{\mathrm{G}}(\mathrm{T}')\) and \(a_g(s)(gtg^{-1}) = gs(t)g^{-1}\) for all \(s \in \mathrm{W}_{\mathrm{G}}(\mathrm{T})\) and all \(t \in T\) .

PROPOSITION 5. a) Every conjugacy class of G meets T.

\begin{enumerate}
\def\labelenumi{\alph{enumi})}
\setcounter{enumi}{1}
\tightlist
\item
  The intersections with \(\mathrm{T}\) of the conjugacy classes of \(\mathbf{G}\) are the orbits of the Weyl group.
\end{enumerate}

Let \(g \in \mathbf{G}\) ; by Th. 2 of no. 2, there exists \(h \in \mathbf{G}\) such that \(g \in h\mathrm{Th}^{-1}\) , hence \(a\) ). By definition of the Weyl group, any two elements in the same orbit of \(\mathrm{W_G(T)}\) on \(T\) are conjugate in \(G\) ; conversely, let \(a, b\) be two elements of \(T\) conjugate under \(G\) . There exists \(h \in G\) such that \(b = hah^{-1}\) ; applying Cor. 7 of Th. 2 (no. 2) with \(A = \{a\}\) , \(s = \operatorname{Int} h\) , \(T' = T\) , we see that there exists \(g \in G\) such that \(\operatorname{Int} hg\) maps \(T\) to \(T\) and \(a\) to \(b\) . The class of \(hg\) in \(\mathrm{W_G(T)}\) then maps \(a\) to \(b\) , hence the proposition.

COROLLARY 1. The canonical injection of T into G defines by passage to the quotient a homeomorphism from \(\mathrm{T}/\mathrm{W}_{\mathrm{G}}(\mathrm{T})\) to the space \(\mathrm{G}/\mathrm{Int}(\mathrm{G})\) of conjugacy classes of G.

Indeed, this is a bijective continuous map between two compact spaces (cf.~General Topology, Chap. III, p.~29, Cor. 1).

COROLLARY 2. Let E be a subset of G stable under inner automorphisms. Then E is open (resp. closed, resp. dense) in G if and only if \(E \cap T\) is open (resp. closed, resp. dense) in T.

This follows from Cor. 1 and the fact that the canonical maps \(T \rightarrow T/W_{G}(T)\) and \(G \rightarrow G/Int(G)\) are open (General Topology, Chap. III, p.~10, Lemma 2).

Denote the Lie algebra of G by g, and that of T by t. The operation of \(W_{\mathrm{G}}(\mathrm{T})\) on T induces a representation, called the canonical representation, of the group \(W_{\mathrm{G}}(\mathrm{T})\) on the R-vector space t.

PROPOSITION 6. a) Every orbit of G on g (for the adjoint representation) meets t.

\begin{enumerate}
\def\labelenumi{\alph{enumi})}
\setcounter{enumi}{1}
\tightlist
\item
  The intersections with \(\mathfrak{t}\) of the orbits of \(\mathrm{G}\) are the orbits of \(\mathrm{W_G(T)}\) on \(\mathfrak{t}\) .
\end{enumerate}

Assertion \(a\) ) follows from Th. 1 (no. 1). Let \(x, y\) be two elements of \(\mathfrak{t}\) conjugate under \(\operatorname{Ad}(\mathrm{G})\) , and let \(h \in \mathrm{G}\) be such that \((\operatorname{Ad} h)(x) = y\) . Applying the corollary of Th. 1 (no. 1) with \(\mathfrak{a} = \{x\}\) , \(u = \operatorname{Ad} h, \mathfrak{t}' = \mathfrak{t}\) , we see that there exists \(g \in \mathrm{G}\) such that \(\operatorname{Ad} hg\) maps \(\mathfrak{t}\) to \(\mathfrak{t}\) and \(x\) to \(y\) . Then \(hg \in \mathrm{N}_{\mathrm{G}}(\mathrm{T})\) (Chap. III, §9, no. 4, Prop. 11), and the class of \(hg\) in \(\mathrm{W}_{\mathrm{G}}(\mathrm{T})\) maps \(x\) to \(y\) , hence the proposition.

COROLLARY. The canonical injection of t into g defines by passage to the quotient a homeomorphism from \(\mathfrak{t}/\mathrm{W}_{\mathrm{G}}(\mathrm{T})\) to \(\mathfrak{g}/\mathrm{Ad}(\mathrm{G})\) .

Denote this map by \(j\) ; it is bijective and continuous (Prop. 6). We have a commutative diagram

\[
\begin{array}{c c c} \mathfrak {t} & \stackrel {{i}} {{\longrightarrow}} & \mathfrak {g} \\ p \Big \downarrow & & \Big \downarrow q \\ \mathfrak {t} / \mathrm{W} _ {\mathrm{G}} (\mathrm{T}) & \stackrel {{j}} {{\longrightarrow}} & \mathfrak {g} / \mathrm{Ad} (\mathrm{G}) \end{array}
\]

where p and q are quotient maps, and i is the canonical injection. Since i and q are proper (General Topology, Chap. I, §10, no. 1, Prop. 2 and General Topology, Chap. III, §4, no. 1, Prop. 2 c)) and since p is surjective, it follows that j is proper (General Topology, Chap. I, §10, no. 1, Prop. 5), and hence is a homeomorphism.

PROPOSITION 7. Let H be a closed subgroup of G containing T.

\begin{enumerate}
\def\labelenumi{\alph{enumi})}
\item
  Denote by \(\mathrm{W}_{\mathrm{H}}(\mathrm{T})\) the subgroup \(\mathrm{N}_{\mathrm{H}}(\mathrm{T})/\mathrm{T}\) of \(\mathrm{W}_{\mathrm{G}}(\mathrm{T})\) ; the group \(H/H_{0}\) is isomorphic to the quotient group \(\mathrm{W}_{\mathrm{H}}(\mathrm{T})/\mathrm{W}_{\mathrm{H}_{0}}(\mathrm{T})\) .
\item
  H is connected if and only if every element of \(\mathrm{W}_{\mathrm{G}}(\mathrm{T})\) that has a representative in H belongs to \(\mathrm{W}_{\mathrm{H}_{0}}(\mathrm{T})\) .
\end{enumerate}

Assertion \(a\) ) follows from Cor. 8 of Th. 2 (no. 2), and assertion \(b\) ) is a particular case of \(a\) ).

